
\chapter{Conclusion}
\label{ch:conclusion}

\section{What was established}

A convolutional classifier trained on curated leaf images reaches a macro-F1 of
0.9931 on held-out curated images and 0.1798 on field images, a difference of 81
points on the same checkpoints. The measurement replicates across five seeds.

The gap is not caused by the model ignoring the plant. Replacing the leaf moves the
prediction to the donor class 50.3\,\% of the time; replacing the background moves it
21.8\,\%, against a 2.6\,\% chance rate and a 4.4\,\% control for the edit itself.
The model uses both, and the leaf carries more of the decision.

Background augmentation during training does not close the gap. Six runs across two
conditions and three seeds give a paired difference of $+1.24$ points with a 95\,\%
interval of $[-3.32, +5.80]$, recovering 1.6\,\% of the gap while costing 0.87 points
in the lab. The direction was predicted from the mechanism measurement before the
experiment ran.

Confidence thresholds do not repair the field case. Calibration error rises by a
factor of 223 between domains, abstaining costs coverage without producing a usable
operating point, and errors are confident errors.

\section{What the project contributes}

The contribution is not a better classifier. It is a measured account of what a
clean-image classifier is worth when the images stop being clean, with the mechanism
tested by intervention rather than asserted, and with the negative result on the
obvious fix pre-registered and explained.

Every number traces to a stored artifact, and the document build fails if a quoted
figure is absent from the registry.

\section{What would be worth doing next}

The results point at leaf appearance as the remaining problem, so the next
experiments are the ones that attack it.

\subsection*{Field-domain leaf augmentation}

Background augmentation failed because it targets the wrong factor. The analogous
treatment for leaf appearance is style transfer restricted to the leaf region, using
field photographs as the style source and PlantVillage leaves as the content. The
harness for regional compositing already exists and the evaluation design in this
report would carry over unchanged.

\subsection*{A small field-labelled fine-tune}

Fine-tuning on a modest number of hand-labelled field images would test how much of
the gap is recoverable appearance versus genuine ambiguity. The number worth knowing
is the slope: how many labelled field images buy how much field macro-F1. If the
curve is steep the problem is practical, and if it is flat the problem is harder than
it looks.

\subsection*{A field-appropriate evaluation set}

PlantDoc at 236 images cannot support intervals tighter than about ten points. A
larger field set, collected with a stated protocol and expert labels for a subset,
would make the remaining questions answerable.

\subsection*{Segment before classifying}

If leaf appearance varies with how the leaf is framed, isolating and normalising the
leaf region before classification may remove variation that no augmentation can
reach. The masks used in the intervention experiment are a starting point, not a
solution.

\section{Closing statement}

A model that reports 99\,\% and delivers 18\,\% is not a model with a bug. It is a
model answering a question nobody asked. The useful work is to find out which
question it is answering, and this report does that for one case, with numbers that
can be checked.
